% !TeX encoding = UTF-8
% !TeX spellcheck = es_ES
% !TeX root = Historia.tex
% !TEX root = Historia.tex
\subsection{Propósito}

Siguiendo la filosofía de tener una guía para desarrollar la maqueta, vamos a exponer
una historia ficticia que justifique los diferentes elementos escénicos y estructurales
de la maqueta, tanto grosso modo, o descripción a alto nivel, como a bajo nivel, pero
que también cuente la historia de la misma y nos permita incluir las diferentes fases. 

El objetivo de esta historia es, en realidad, el contexto escénico, por lo que debemos
tenerlo lo antes posible. Aun así debemos recordar que es un cuento, a menos que queramos
recrear una línea ferroviaria exacta, y va a ser muy difícil que podamos escribir toda la
historia antes de empezar a diseñar nuestra maqueta. Por lo tanto este cuento lo iremos
escribiendo poco a poco. 
 
No lo confundamos con el objetivo en si mismo de la maqueta. Con la maqueta, además de
divertirnos, podemos querer hacer algo más. Pero la justificación de porque esta ese
elemento en ese sitio concreto, es el objetivo de esta historia o contexto. Léase porque
hay una industria papelera, o una estación perdida en la montaña, o… Si bien es cierto
que el objetivo de una maqueta puede ser la exposición de un diorama de una historia
concreta, con lo que ambos objetivos, en este caso, coincidirían. 
 
Los maquetistas noveles pueden verse abrumados en pensar esta historia, cuando lo que
quieren normalmente es tener un óvalo para rodar sus modelos. Pero tener esta historia
nos permitirá tener claro las intenciones, saber que modelos comprar y como diseñar
la maqueta.  
 
Esta historia conviene que sea un documento vivo, que se vaya modificando y creciendo
según lo que podamos hacer en cada momento y un poquito más, cuanto más completa sea,
más podremos avanzar en nuestra mente y papel de como queramos que sea nuestra maqueta
en el futuro. Y cuanto más detalle, más fácil será diseñar y más barato. 
 
Entonces, lo que debemos tener lo antes posible es el tema, leitmotiv, o motivación,
que deberá ser resumen de un párrafo, de esta historia.  No es necesario sentarse al
principio del diseño para crear esta historia y así luego pensar en la maqueta. Pero si
conviene pararse a pensar en que Objetivos se tiene con la maqueta, que queremos hacer
y definir una historia que lo justifique. 

\subsection{Requisitos}
El punto de partida para una maqueta, debe ser la respuesta a ¿Qué maqueta quiero hacer?.
Como hemos dicho en el capítulo correspondiente son los requisitos. Los tenemos que tener
claro antes de hacer este documento. 
 
Así mismo este es buen momento para pensar también que nos gustaría, además tener, e
incluirlos como requisitos de baja prioridad. Finalmente veremos el material que tenemos
para ver como integrarlo y justificarlo en la maqueta. 
 
Recordemos que para Daniel Bahn queremos tener, descartando los más técnicos:
\begin{itemize}
\item Representar tres escenas inspiradas en sitios reales, por importancia sentimental:
\begin{itemize}
\item La estación del puerto de La Coruña hasta la playa de lazareto 
\item La estación del Santo Sepulcro de Zaragoza 
\item Una estación de montaña como la de Canfranc 
\end{itemize}
\item Tener un una vía continua para realizar fotos, sin preocupaciones de tener
que evitar choques contra fin de vía 
\item Tener una zona de puzzle de clasificación 
\item Ser más o menos realista en trazado y operación. 
\end{itemize}
 
Así mismo queremos mantener el material existente, siendo la mayoría material alemán de
épocas modernas (IV - V) con algún vagón III. También se quiere establecer como fechas 
los años finales de los 90 y principios del 2000, antes del boom de la alta velocidad
y como las escenas que queremos representar son de esa época, pues se justifica. 

\subsection{Desarrollo de la Historia}
Para desarrollar la historia, ficticia, de la maqueta empezaremos por buscar los
requisitos que más impactan en la parte histórica. En un primer vistazo, y sin fijarnos
mucho en las escenas, lo primero que podemos fijarnos que se quiere representar escenas
Españolas con material Alemán. 
 
En segundo lugar, y pensando ya en las escenas, veremos que el puerto de La Coruña es
relativamente moderno (años 60 aprox.) y la estación de Zaragoza cerro en esa época por
una más moderna. Canfranc ha estado desde los años 20, y es ahora cuando se está
modificando sustancialmente, aun así el requisito es una estación de montaña. 
 
Esto son los requisitos que más pueden chocar entre sí por tener que representar lugares
y tiempos dispares. El resto si se puede justificar con una historia mejor, pero siempre
se pueden solucionar sin ello. 
 
El tema de material alemán en paisajes españoles se puede solucionar de muchas maneras,
como por ejemplo que el sector está liberado y cualquier compañía puede operar y en este
caso es Daniel Bahn la única que lo hace. Otra opción es que el material alemán está
barato por cualquier razón. Pero para que esto esté justificado, nos aprovechamos de
la historia real, modificándola en una época anterior a la deseada y acabar en una
ligeramente posterior a la misma. Entonces en Daniel Bahn cambiaremos ligeramente el
resultado de las guerras mundiales cambiando ligeramente las consecuencias de las mismas,
donde a los alemanes se le exige más pago y eso produce una expansión de su tecnología
por toda Europa y como queremos fijarnos en principios de los 90, podemos acabar cuando
llega la Alta Velocidad en Zaragoza, coincidiendo con la exposición universal de Zaragoza.